
\documentclass{article}
\usepackage{colortbl}
\usepackage{makecell}
\usepackage{multirow}
\usepackage{supertabular}

\begin{document}

\newcounter{utterance}

\centering \large Interaction Transcript for game `cladder', experiment `full\_v1.5\_default', episode 8797 with reflexion{-}Qwen3.5{-}4B.
\vspace{24pt}

{ \footnotesize  \setcounter{utterance}{1}
\setlength{\tabcolsep}{0pt}
\begin{supertabular}{c@{$\;$}|p{.15\linewidth}@{}p{.15\linewidth}p{.15\linewidth}p{.15\linewidth}p{.15\linewidth}p{.15\linewidth}}
    \# & \multicolumn{2}{c}{Player} && \multicolumn{2}{c}{Game Master} \\
    \hline

    \theutterance \stepcounter{utterance}  
    & & & \multicolumn{4}{p{0.6\linewidth}}{
        \cellcolor[rgb]{0.9,0.9,0.9}{
            \makecell[{{p{\linewidth}}}]{
                \texttt{\tiny{[P1$\langle$GM]}}
                \texttt{You have attempted this same task before and did not succeed. These are the plans you wrote after those attempts.} \\
\texttt{Trial \#1: New plan: I failed to complete the task because I stopped providing the required reasoning after the first bullet point, leaving the answer incomplete. The prompt explicitly requires a full explanation supporting the "Yes" or "No" decision, including the calculation for the counterfactual group (\$Y\_\{yupt\}=0\$). My strategy was to calculate the probability for the treated group (\$Y\_\{yupt\}=1\$) correctly but neglected to calculate the probability for the control group (\$Y\_\{yupt\}=0\$) using the provided conditional probabilities (\$P(zupj=1 \textbackslash{}mid yupt=0) = 51\textbackslash{}\%\$ and \$P(zupj=0 \textbackslash{}mid yupt=0) = 49\textbackslash{}\%\$) and the corresponding outcome probabilities. Next time, I will first determine the marginal probability of \$muvq\$ for both groups (\$Y\_\{yupt\}=1\$ and \$Y\_\{yupt\}=0\$) by weighting the conditional probabilities with the conditional probabilities of \$zupj\$, and then compare these two values to definitively answer the counterfactual question.} \\
\\ 
\\ 
\texttt{You are an expert in causal inference. The following question is not a typical commonsense query, but rather a meticulously designed question created by a professor specializing in causal inference, intended to assess the students' mastery of the course content.} \\
\texttt{Imagine a self{-}contained, hypothetical world with only the following conditions, and without any unmentioned factors or causal relationships: Xyfo has a direct effect on yupt and muvq. Yupt has a direct effect on zupj. Zupj has a direct effect on muvq. Xyfo is unobserved. For those who are not yupt and are not zupj, the probability of muvq is 90\%. For those who are not yupt and are zupj, the probability of muvq is 54\%. For those who are yupt and are not zupj, the probability of muvq is 96\%. For those who are yupt and are zupj, the probability of muvq is 57\%. For those who are not yupt, the probability of zupj is 51\%. For those who are yupt, the probability of zupj is 34\%. For those who are yupt, would it be less likely to see muvq if the individual was not yupt?} \\
\texttt{Start your answer with 'Yes' or 'No', followed by additional reasoning or evidence to support your explanation.} \\
            }
        }
    }
     \\ \\

    \theutterance \stepcounter{utterance}  
    & \multicolumn{4}{p{0.6\linewidth}}{
        \cellcolor[rgb]{0.561,0.737,0.561}{
            \makecell[{{p{\linewidth}}}]{
                \texttt{\tiny{[P1$\rangle$GM]}}
                \texttt{Yes} \\
\\ 
\texttt{To determine whether it would be less likely to see \$muvq\$ if an individual treated as \$yupt=1\$ were actually in the counterfactual state of \$yupt=0\$, we must compare the marginal probability of \$muvq\$ under the treated condition (\$P(muvq \textbackslash{}mid yupt=1)\$) with the marginal probability of \$muvq\$ under the counterfactual control condition (\$P(muvq \textbackslash{}mid yupt=0)\$).} \\
\\ 
\texttt{**Step 1: Calculate the marginal probability of \$muvq\$ for the treated group (\$yupt=1\$).**} \\
\texttt{The problem provides the conditional probabilities of \$muvq\$ based on the status of \$zupj\$ for those who are \$yupt\$:} \\
\texttt{*   \$P(muvq \textbackslash{}mid yupt=1, zupj=1) = 57\textbackslash{}\% = 0.57\$} \\
\texttt{*   \$P(muvq \textbackslash{}mid yupt=1, zupj=0) = 96\textbackslash{}\% = 0.96\$} \\
\\ 
\texttt{We are also given the conditional probabilities of \$zupj\$ for those who are \$yupt\$:} \\
\texttt{*   \$P(zupj=1 \textbackslash{}mid yupt=1) = 34\textbackslash{}\% = 0.34\$} \\
\texttt{*   \$P(zupj=0 \textbackslash{}mid yupt=1) = 1 {-} 0.34 = 66\textbackslash{}\% = 0.66\$} \\
\\ 
\texttt{Using the law of total probability, we calculate the overall probability for the treated group:} \\
\texttt{\$\$P(muvq \textbackslash{}mid yupt=1) = P(muvq \textbackslash{}mid yupt=1, zupj=1) \textbackslash{}cdot P(zupj=1 \textbackslash{}mid yupt=1) + P(muvq \textbackslash{}mid yupt=1, zupj=0) \textbackslash{}cdot P(zupj=0 \textbackslash{}mid yupt=1)\$\$} \\
\texttt{\$\$P(muvq \textbackslash{}mid yupt=1) = (0.57 \textbackslash{}times 0.34) + (0.96 \textbackslash{}times 0.66)\$\$} \\
\texttt{\$\$P(muvq \textbackslash{}mid yupt=1) = 0.1938 + 0.6336 = 0.8274 \textbackslash{}text\{ (or \} 82.74\textbackslash{}\%)\$\$} \\
\\ 
\texttt{**Step 2: Calculate the marginal probability of \$muvq\$ for the counterfactual group (\$yupt=0\$).**} \\
\texttt{The problem provides the conditional probabilities of \$muvq\$ based on the status of \$zupj\$ for those who are *not* \$yupt\$:} \\
\texttt{*   \$P(muvq \textbackslash{}mid yupt=0, zupj=1) = 54\textbackslash{}} \\
            }
        }
    }
    & & \\ \\

    \theutterance \stepcounter{utterance}  
    & & & \multicolumn{2}{p{0.3\linewidth}}{
        \cellcolor[rgb]{0.9,0.9,0.9}{
            \makecell[{{p{\linewidth}}}]{
                \texttt{\tiny{[GM$|$GM]}}
                \texttt{Yes} \\
            }
        }
    }
    & & \\ \\

    \theutterance \stepcounter{utterance}  
    & & & \multicolumn{2}{p{0.3\linewidth}}{
        \cellcolor[rgb]{0.9,0.9,0.9}{
            \makecell[{{p{\linewidth}}}]{
                \texttt{\tiny{[GM$|$GM]}}
                \texttt{yes} \\
            }
        }
    }
    & & \\ \\

    \theutterance \stepcounter{utterance}  
    & & & \multicolumn{2}{p{0.3\linewidth}}{
        \cellcolor[rgb]{0.9,0.9,0.9}{
            \makecell[{{p{\linewidth}}}]{
                \texttt{\tiny{[GM$|$GM]}}
                \texttt{game\_result = WIN} \\
            }
        }
    }
    & & \\ \\

\end{supertabular}
}

\end{document}
